% Основной текст статьи. Подключается из article.tex шаблона МИСИС.
% Блоки названия, авторов, аннотации и ключевых слов задаются в article.tex
% в соответствии с командами шаблона.

\section{Введение}

Электроэнцефалография (ЭЭГ) используется в интерфейсах «мозг--компьютер» для распознавания намерений пользователя по электрической активности мозга. Одной из распространённых парадигм является моторное воображение: участник представляет движение без его фактического выполнения. В наборе BCI Competition IV 2a представлены четыре класса моторного воображения: левая рука, правая рука, ноги и язык; записи включают 22 канала ЭЭГ и получены с частотой дискретизации 250 Гц \cite{Tangermann2012,BCICompetitionIV}.

Большинство классических подходов опирается на спектральные характеристики, пространственные фильтры или их комбинации. При этом ЭЭГ является нестационарным сигналом, а её структура может содержать нелинейные и масштабно-зависимые особенности. Фрактальные характеристики позволяют описывать некоторые свойства сложности и нерегулярности сигнала. Однако наличие фрактального признака само по себе не означает, что он улучшает классификацию: это необходимо проверять экспериментально на едином протоколе и с независимой оценкой качества.

В настоящей работе рассматриваются две характеристики, основанные на разных подходах: фрактальная размерность Хигучи (Higuchi fractal dimension, HFD) и корреляционная размерность $D_2$, оцениваемая после реконструкции пространства состояний. Метод Хигучи предназначен для оценки фрактальной размерности графика временного ряда \cite{Higuchi1988}. Подход Грассбергера--Прокаччиа оценивает масштабный показатель по зависимости корреляционной суммы от радиуса \cite{GrassbergerProcaccia1983}. Эти характеристики исследуются как потенциальные признаки классификации моторного воображения.

Цель работы — исследовать применимость фрактальных характеристик ЭЭГ-сигналов для классификации состояний моторного воображения и оценить устойчивость получаемых признаков к параметрам расчёта. Для достижения цели необходимо: (1) реализовать расчёт HFD и $D_2$; (2) проверить реализацию на синтетических сигналах с известными качественными свойствами; (3) исследовать чувствительность оценки $D_2$ к задержке, размерности вложения и диапазону радиусов; (4) сравнить качество классификации по фрактальным признакам с базовыми признаками ЭЭГ при едином протоколе валидации.

\section{Материалы и методы}

\subsection{Набор данных}

В качестве источника данных используется BCI Competition IV, набор 2a. Он содержит записи девяти участников, четыре класса моторного воображения, 22 EEG-канала и три EOG-канала; частота дискретизации составляет 250 Гц \cite{Tangermann2012,BCICompetitionIV}. EOG-каналы не включаются в набор признаков классификатора. В пилотных экспериментах использована тренировочная запись участника A01 (файл A01T.gdf). Данная запись предназначена для отладки и анализа чувствительности метода и не рассматривается как достаточная выборка для общих выводов о классификации.

Для основного эксперимента предполагается использовать тренировочные записи всех доступных участников. Сигналы должны быть разделены на эпохи относительно маркеров начала моторного воображения. В пилотном анализе выделялись эпохи длительностью 4 с после маркера, использовался диапазон 8--30 Гц и канал C3. Этот выбор следует понимать как диагностический, а не как окончательно обоснованный оптимальный протокол. Перед итоговым сравнением параметры фильтрации, временное окно и набор каналов фиксируются для всех сравниваемых методов.

\subsection{Фрактальная размерность Хигучи}

Метод Хигучи строит последовательности, получаемые из исходного ряда с различными шагами, и анализирует зависимость длины этих последовательностей от масштаба. Наклон линейного участка в логарифмических координатах используется как оценка HFD \cite{Higuchi1988}. В рамках эксперимента HFD рассчитывается отдельно для каналов ЭЭГ; параметры алгоритма, прежде всего максимальный шаг $k_{\max}$, фиксируются до сравнения классификаторов.

HFD является характеристикой графика исходного временного ряда. Она не равна корреляционной размерности реконструированного аттрактора и не должна интерпретироваться как прямая замена $D_2$.

\subsection{Реконструкция пространства состояний и корреляционная размерность}

Для одномерного временного ряда $x(t)$ формируется задержанное вложение:
\begin{equation}
\mathbf{x}_i^{(m,\tau)} = [x_i, x_{i+\tau}, \ldots, x_{i+(m-1)\tau}],
\end{equation}
где $m$ — размерность вложения, а $\tau$ — задержка в отсчётах. Для реконструированных векторов вычисляется корреляционная сумма
\begin{equation}
C_m(r)=\frac{2}{N_r(N_r-1)}\sum_{i<j,\,|i-j|>W}
\mathbb{I}\bigl(\|\mathbf{x}_i-\mathbf{x}_j\|<r\bigr),
\end{equation}
где $r$ — радиус, $W$ — окно Тейлера, уменьшающее влияние близких по времени точек, а $\mathbb{I}$ — индикатор выполнения условия. Обозначение $N_r$ здесь относится к числу допустимых векторов после учёта граничных условий и окна Тейлера; конкретная реализация должна использовать одинаковое определение для всех экспериментов.

Корреляционная размерность оценивается по наклону зависимости $\log C_m(r)$ от $\log r$ на выбранном диапазоне масштабов:
\begin{equation}
D_2 \approx \frac{d\log C_m(r)}{d\log r}.
\end{equation}
Для интерпретации оценки необходимо наличие устойчивого участка масштабирования. Высокий коэффициент детерминации линейной регрессии на коротком участке сам по себе не доказывает существование такого участка. Поэтому анализ включает чувствительность к диапазону радиусов, задержке $\tau$, размерности вложения $m$ и окну Тейлера. Если устойчивый участок не удаётся обосновать, результат рассматривается как диагностический наклон, а не как подтверждённая оценка $D_2$.

\subsection{Проверка на синтетических сигналах}

Реализация проверяется на синусоидальном сигнале, зашумлённой синусоиде, сигнале с двумя частотными компонентами и белом шуме. Цель проверки — выявить очевидные ошибки реализации и продемонстрировать, что оценка зависит от структуры сигнала и выбранного диапазона масштабов. Синтетические сигналы не заменяют проверку на ЭЭГ и не служат доказательством эффективности признаков для классификации моторного воображения.

\subsection{Протокол классификации}

Для итогового эксперимента сравниваются по меньшей мере следующие наборы признаков: (1) спектральные признаки базового уровня; (2) HFD; (3) признаки, основанные на $D_2$, только в тех условиях, где диапазон масштабирования методологически обоснован; (4) объединённые спектральные и фрактальные признаки. Один и тот же классификатор и одинаковый протокол валидации используются для сопоставимых наборов признаков. В качестве одного из базовых классификаторов может использоваться SVM с RBF-ядром и балансировкой классов.

Чтобы оценить переносимость между участниками, основным вариантом разделения является проверка с исключением одного участника: все эпохи одного участника остаются в тестовой части, а обучение выполняется на остальных участниках. Любая стандартизация признаков подгоняется только на обучающей части. Основные метрики — сбалансированная точность (balanced accuracy), macro-F1 и матрица ошибок. Сравнение проводится на одних и тех же разбиениях. Для утверждений об улучшении дополнительно требуется парное статистическое сравнение и поправка на множественные проверки, если анализируется несколько каналов или вариантов параметров.

\section{Результаты предварительных экспериментов}

\subsection{Синтетические сигналы}

Предварительная проверка на синтетических сигналах показала, что результат зависит от диапазона радиусов, используемого для регрессии. Для чистой синусоиды на подходящих диапазонах наблюдались наклоны, близкие к единице. Для сигнала с двумя частотными компонентами на части параметров получались значения около двух, однако при увеличении размерности вложения оценки менялись. Для белого шума наклон возрастал с размерностью вложения; при этом коэффициент детерминации мог оставаться высоким. Это подтверждает, что одного высокого $R^2$ недостаточно для выбора диапазона масштабирования.

\subsection{Пилотный анализ ЭЭГ}

На первой тренировочной записи A01T.gdf исследовался канал C3 после полосовой фильтрации 8--30 Гц. В одном диагностическом эпизоде левой руки при задержке $\tau=4$ отсчёта медианные наклоны возрастали с размерностью вложения: приблизительно 1.90 при $m=2$, 2.64 при $m=3$, 3.14 при $m=4$, 3.62 при $m=5$, 4.10 при $m=6$, 4.89 при $m=8$ и 5.65 при $m=10$. Эти значения не демонстрируют очевидного плато по размерности вложения. Кроме того, изменение диапазона радиусов приводило к заметному изменению оценок.

В отдельном пилотном сравнении нескольких эпох левой и правой руки для одного участника устойчивое плато также не было однозначно установлено; различия между классами были небольшими и меняли направление в зависимости от параметров. Поскольку анализ ограничен одним участником и небольшим числом эпох, он является только диагностическим. По этим данным нельзя делать вывод о статистически значимом различии классов или о качестве классификации на популяционном уровне.

\subsection{Промежуточная интерпретация}

Полученные наблюдения указывают прежде всего на необходимость контролировать устойчивость оценки $D_2$. Если диапазон масштабирования не удаётся выбрать воспроизводимо, принудительное сведение данных к одному числу может вносить нестабильность в признаки классификатора. Для HFD и других фрактальных характеристик также необходима проверка вклада признаков в рамках сопоставимого разделения обучающей и тестовой выборок. До завершения этой проверки нельзя утверждать, что добавление фрактальных характеристик повышает качество классификации.

\section{Обсуждение}

Работа рассматривает два отличающихся способа описания нерегулярности ЭЭГ: оценку фрактальной размерности графика временного ряда и анализ масштабного поведения корреляционной суммы реконструированного пространства. Их сравнительная ценность определяется не только различиями средних значений между классами, но и устойчивостью к параметрам метода, межиндивидуальной вариативностью и итоговой метрикой классификации.

Ограничением текущего этапа является то, что расчёты $D_2$ пока носят пилотный характер: анализируются отдельные каналы и ограниченная часть данных, а устойчивый диапазон масштабирования не подтверждён для всех рассматриваемых условий. Возможные источники нестабильности включают конечную длину эпох, шум, фильтрацию, выбор задержки, размерности вложения, радиусов и окна Тейлера. Дополнительное ограничение — невозможность переносить выводы с одного участника на других без межсубъектной проверки.

Следующий этап должен включать фиксированный протокол предобработки для всех девяти участников, анализ HFD на выбранном диапазоне параметров, строгую проверку устойчивости $D_2$, а затем сравнение классификаторов с фрактальными и базовыми признаками. Данные тестовых участников нельзя использовать для настройки параметров или стандартизации.

\section{Заключение}

Подготовлена экспериментальная схема исследования фрактальных характеристик ЭЭГ для классификации моторного воображения, включающая HFD, реконструкцию пространства состояний и оценку корреляционной размерности $D_2$. Проверка на синтетических сигналах и пилотный анализ одной записи BCI Competition IV 2a выявили чувствительность оценок корреляционного наклона к выбору диапазона радиусов и размерности вложения. В доступном пилотном анализе устойчивое плато $D_2$ не было убедительно установлено, поэтому эти оценки пока нельзя считать надёжными признаками для окончательного классификационного сравнения.

Дальнейшая работа направлена на воспроизводимое извлечение признаков по набору данных, межсубъектную оценку качества классификации и прямое сравнение фрактальных признаков со спектральным базовым уровнем. Итоговый вывод о полезности фрактальных характеристик будет сделан только после проведения этого сравнения.
